\section{Uniform bounds for quartic dimension fields}\label{sec:quartic-bound}

The following bound applies to arbitrary integral orders in the
dimension field. Its cutoff $950$ and resulting multiplicity $30$
are auxiliary refinements; the final finite proof uses the common
cutoff $1000$ and multiplicity $31$.
\begin{theorem}\label{thm:quartic-uniform-bound}
Let $\C$ be a self-dual rank-four complex fusion category, and
suppose that its Frobenius--Perron dimension $D$ generates a quartic
field $K$. If $D\ge950$, then $\C$ is spherical.
\end{theorem}

We prove the theorem in two field cases. First suppose
$K\ne\Q(\sqrt2,\sqrt5)$, and assume for a contradiction that
$\C$ is nonspherical with $D\ge950$. We use the sphericalization notation $R,T,D,E$ and the
central-column identities. By
Corollary~\ref{col:cor:rank-four-dimension-field-descent}, both rational
algebras are the same totally real abelian quartic field. Choose the
largest conjugate $E$ of the signed codegree and put
\[
 \alpha=E/D,\qquad p_\sigma=1/\sigma(D),\qquad
 q_\sigma=1/\sigma(E),\qquad d=p_0=1/D.
\]
Then $\sum p_\sigma=\sum q_\sigma=1$, and $d\le q_0$ are their
respective least entries. Central integrality and the order-discriminant
identity show that $\alpha$ is a totally positive integral $d$-number
whose norm $h$ is a positive integer square. In particular,
$\prod q_\sigma\le\prod p_\sigma$. Moreover $\alpha_0$ is the least
conjugate: the signed central constituent labelled by $\sigma$ has
ordinary folded dimension $D/\sigma(D)$ and signed dimension
$E/\sigma(E)$, so spectral domination gives
\[
 E/\sigma(E)\le D/\sigma(D),\qquad
 E/D\le\sigma(E)/\sigma(D).
\]

Write $u=\sum p_\sigma^2$ and $v=\sum q_\sigma^2$.
Corollary~\ref{core:quartic-short-moments} gives $u+v\le4/5$ if
$\sqrt2\notin K$, and $2u+v\le9/8$ if $\sqrt5\notin K$.
In the latter case
\[
 u\ge d^2+\frac{(1-d)^2}{3}>\frac{13}{40}
 \qquad(D\ge100)
\]
also gives $u+v<4/5$. A quartic field containing both square roots
is their compositum. Thus both branches reduce to
\begin{equation}\label{eq:quartic-two-moments}
 u+v\le\frac45.
\end{equation}
Schur's inequality and signed trace domination give
\begin{equation}\label{eq:quartic-norm-estimates}
 v\ge\frac12-\frac9{2h}e_3(p),\qquad
 e_3(p)\le\frac{(1-d)^2(1+8d)}{27}.
\end{equation}
If $h\ge9$, these imply
\[
 u+v\ge\frac12+\frac{(1-d)^2(17-8d)}{54}>\frac45
 \qquad(D\ge100),
\]
a contradiction. Therefore
\begin{equation}\label{eq:quartic-small-norm}
 h\in\{1,4\}.
\end{equation}

\subsection{Product and ratio estimates}

\begin{lemma}\label{lem:three-largest-product}
Let positive quadruples $p,q$ each have coordinate sum one. Suppose
$p_0=d\le1/950$, $q_0=\min q_i$, $u+v\le4/5$, and
$\prod q_i\le\prod p_i$, where $u=\sum p_i^2$, $v=\sum q_i^2$.
Then $q_0<1/168$.
\end{lemma}
\begin{proof}
Put
\[
 d_*=\frac1{950},\quad r=\frac1{168},\quad a=\frac7{250},
 \quad S=1-r,\quad
 c=\frac45-d_*^2-\frac{(1-d_*)^2}{3}.
\]
By monotonicity for $0<d\le d_*$,
\[
 v\le c,\qquad
 \prod q_i\le\prod p_i\le\frac{d(1-d)^3}{27}
 \le\frac{d_*(1-d_*)^3}{27}.
\]
Suppose all $q_i\ge r$, and minimize their product subject to this
lower bound, their sum being one, and their square sum being at most
$c$. A minimizer has a coordinate equal to $r$. Indeed, an interior
minimizer with inactive square-sum constraint would have all coordinates
equal, which is a local maximum. With the constraint active, the
multiplier equations for the logarithm of the product give at most
two values $a_0<b_0$. The constrained Hessian is diagonal with entries
$1/(a_0b_0)-1/q_i^2$. If the smaller value occurs at least twice, the tangent
vector increasing one such coordinate and decreasing another has
negative second variation. If it occurs only once, the square sum is
$a_0^2+(1-a_0)^2/3<1/3<c$, again impossible. The two constraint
gradients are independent at every such two-value point.

At a boundary coordinate $r$, the other three have sum $S$ and square
sum at most $c-r^2<S^2/2$. Their unconstrained positive product minimum
has one small and two equal large coordinates,
\[
 \left(t,\frac{S-t}{2},\frac{S-t}{2}\right),\qquad
 t^2+\frac{(S-t)^2}{2}=c-r^2,\quad 0<t<S/3.
\]
This is the usual two-value multiplier calculation for three variables;
the one-small-coordinate branch has the smaller product.
In fact $t>a>r$, because
\[
 a^2+\frac{(S-a)^2}{2}-(c-r^2)
 =\frac{2706003}{35378000000}>0
\]
and the quadratic decreases up to $S/3$. Since $x(S-x)^2$ increases
on that interval, every admissible quadruple consequently has product
greater than $ra(S-a)^2/4$. But
\[
 \frac{ra(S-a)^2}{4}-\frac{d_*(1-d_*)^3}{27}
 =\frac{30295201169}{1379317464000000000}>0,
\]
contradicting the preceding upper bound.
\end{proof}

% Exact ratio lemmas and normalized-moment argument. No preamble.
\begin{lemma}\label{lem:triple-ratios-950}
Let $P,Q$ be positive triples whose coordinates sum to one, and put
$U=\sum P_i^2$, $V=\sum Q_i^2$. Then
\[
 U+V\le\frac{201}{250}\ \Longrightarrow\ 
 \sum_i\frac{P_i}{Q_i}<\frac{51}{4},\qquad
 U+V\le\frac{8063}{10000}\ \Longrightarrow\ 
 \sum_i\frac{P_i}{Q_i}<\frac{68}{5}.
\]
\end{lemma}
\begin{proof}
Put $r=\sum Q_i^{-1}$. Cauchy--Schwarz in the plane of coordinate sum
zero gives
\[
 \sum_i\frac{P_i}{Q_i}\le
 \frac r3+\sqrt{\left(B-V-\frac13\right)
             \left(\frac{2r^2}{3}-\frac{4r}{1-V}\right)}.
\]
Here $B$ denotes the relevant bound on $U+V$. For fixed $V<1/2$ the
right-hand side increases with $r\ge9$. The product-minimizing positive
triple with prescribed sum and square sum has one small and two equal
large coordinates. Consequently it suffices to take
\[
 Q=\left(t,\frac{1-t}{2},\frac{1-t}{2}\right),\qquad
 P=\left(a,\frac{1-a}{2},\frac{1-a}{2}\right),
\]
where
\[
 a=\frac13+\sqrt{\frac{2(B-2/3)}3-(t-1/3)^2},\qquad
 \frac{1-\sqrt{6B-4}}3\le t\le\frac13.
\]
For the proposed ratio bound $R$, define
\[
 N(t)=-3Rt^2+(3R-3)t-1,\qquad
 F(t)=N(t)^2+(1-3t)^4-(6B-4)(1-3t)^2.
\]
The required inequality follows from $N(t)>0$ and $F(t)>0$ by
multiplying positive denominators and squaring positive sides; the
endpoint $t=1/3$ gives ratio $3$ directly. Indeed,
\[
 \frac at+\frac{2(1-a)}{1-t}<R
 \quad\Longleftrightarrow\quad
 a-\frac13<\frac{N(t)}{3(1-3t)}
 \qquad(t<1/3).
\]

For $(B,R)=(201/250,51/4)$, one has $t>3/100$,
$N(3/100)=923/40000>0$, and $N$ is increasing up to $1/3$.
The polynomial $f=(2000/3)F$ is
\[
 f(t)=1029375t^4-1869750t^3+910431t^2-51704t+784.
\]
Writing $x=t-1/32$ gives the exact identity
\[
\begin{split}
 f(t)&=H(t)x^2+Lx+C,\\
 H(t)&=1029375t^2-\frac{28886625}{16}t
                         +\frac{815705469}{1024},\\
 L&=-\frac{1263121}{8192},\qquad C=\frac{1325631}{1048576}.
\end{split}
\]
The function $H$ decreases on $[0,1/3]$, and
$H(1/3)=316577469/1024>300000$, $|L|<155$, $C>1$.
Therefore
\[
 f(t)>300000x^2-155|x|+1>0,
\]
since $155^2<1200000$.

For $(B,R)=(8063/10000,68/5)$, one has $t>1/36$,
$N(1/36)=1/54>0$, and $N$ is again increasing up to $1/3$.
Now $f=(5000/3)F$ equals
\[
 f(t)=2909400t^4-5320800t^3+2594833t^2-137622t+1937.
\]
Writing $x=t-29/1000$ gives
\[
\begin{split}
 f(t)&=H(t)x^2+Lx+C,\\
 H(t)&=2909400t^2-\frac{25760274}{5}t
                            +\frac{11467835081}{5000},\\
 L&=-\frac{327793717}{1250000},\qquad
 C=\frac{2526625707}{5000000000}.
\end{split}
\]
Again $H$ decreases on $[0,1/3]$, and
$H(1/3)=13492231243/15000>800000$, $|L|<263$, $C>1/2$.
Thus
\[
 f(t)>800000x^2-263|x|+\frac12>0,
\]
because $263^2<1600000$.
\end{proof}

\begin{lemma}\label{lem:relative-trace-950}
Let positive quadruples $p,q$ each have coordinate sum one, write
$d=p_0$, $q=q_0$, $u=\sum p_i^2$, $v=\sum q_i^2$, and suppose
\[
 0<d\le\frac1{950},\qquad d\le q<\frac1{168},\qquad
 u+v\le\frac45.
\]
Then $\sum_{i=0}^3p_i/q_i<14$.
\end{lemma}
\begin{proof}
Normalize the last three coordinates by
$P_i=p_i/(1-d)$, $Q_i=q_i/(1-q)$, and write
$U=\sum P_i^2$, $V=\sum Q_i^2$. Since $U\ge1/3$ and $q\ge d$,
\[
 U+V\le B(d,q):=\frac13+
 \frac{4/5-d^2-(1-d)^2/3-q^2}{(1-q)^2}.
\]
The function $B$ increases in both variables on the indicated domain.
If $q\le1/300$, then
\[
 B(d,q)\le B(1/950,1/300)
       =\frac{77826934}{96821283}<\frac{201}{250}.
\]
Lemma~\ref{lem:triple-ratios-950} gives
\[
 \sum_{i=0}^3\frac{p_i}{q_i}
 <1+\frac{300}{299}\frac{51}{4}
 =14-\frac{62}{299}<14.
\]
If $q>1/300$, then
\[
 B(d,q)<B(1/950,1/168)
       =\frac{15220412926}{18877366875}<\frac{8063}{10000},
\]
and the other ratio bound gives
\[
 \sum_{i=0}^3\frac{p_i}{q_i}
 <\frac{300}{950}+\frac{168}{167}\frac{68}{5}
 =14-\frac{44}{15865}<14.
\]
\end{proof}

\subsection{A finite arithmetic lemma}

\begin{lemma}\label{lem:abelian-low-trace}
Let $\alpha$ be a totally positive integral $d$-number of degree four,
with abelian Galois field, norm $h\in\{1,4\}$ and trace at most $13$.
Its minimal polynomial $x^4-tx^3+sx^2-rx+h$ is one of the following.
The column $z$ gives an integer root of the cubic resolvent.
\[
\begin{array}{rrrrrl|rrrrrl}
t&s&r&h&z&G&t&s&r&h&z&G\\\hline
7&14&8&1&3&C_4&12&24&12&1&2&V_4\\
8&14&7&1&3&C_4&12&28&12&1&2&V_4\\
8&19&12&1&3&C_4&12&30&24&4&4&V_4\\
8&20&16&1&4&V_4&12&32&24&4&4&V_4\\
9&19&9&1&2&C_4&12&34&20&1&6&C_4\\
9&26&24&1&7&C_4&12&34&23&1&3&C_4\\
10&26&20&4&4&V_4&12&34&28&1&18&C_4\\
11&21&11&1&2&C_4&12&45&54&1&9&V_4\\
11&31&11&1&2&C_4&13&29&17&1&18&C_4\\
12&19&8&1&3&C_4&13&40&13&1&2&C_4\\
12&23&12&1&2&V_4&13&49&57&1&18&C_4\\
\end{array}

\]
If $g(x)=(x+2)^2/(x^2+1)$ and $\alpha_0$ is the least conjugate,
then
\begin{equation}\label{eq:low-trace-scalar}
 \sum_{i=1}^3g(\alpha_i)<\frac{423}{40}.
\end{equation}
The same inequality holds for a nonrational quadratic $d$-number whose
four-dimensional norm is $1$ or $4$ and whose four-dimensional trace is
at most $13$, counting each conjugate twice.
\end{lemma}
\begin{proof}
Here are complete integer tests for the finite table. An abelian
quartic group fixes a partition of the four positive roots into pairs.
Write their sums as $a,a'$ and products as $b,b'$, and put $z=b+b'$.
Then $z$ is an integer and
\[
 A=(a-a')^2=t^2-4s+4z,\quad B=(b-b')^2=z^2-4h,\quad
 AB=(tz-2r)^2.
\]
Positivity and symmetric means give
\[
 4\le t\le13,\quad 2\sqrt h\le z\le t^2/4-2\sqrt h,
 \quad z<s\le3t^2/8,\quad 1\le r\le t^3/16.
\]
For $h=4$ the $d$-number criterion further requires $2\mid t,s$ and
$4\mid r$. For each $h,t,z,s$ in these bounds, retain $A\ge0$ with
$AB$ a square and set $r=(tz\pm\sqrt{AB})/2$.
Put
\[
 S=t^2-2s-2z,\qquad J=s^2+2sz-3z^2-4tr+16h.
\]
The two pair discriminants have sum $S$ and product $J$; therefore
all four roots are positive precisely when $S,J>0$, after reducible
polynomials are discarded. Reducibility is tested by positive divisors
of $h$ for a linear factor, or positive integers $b,d$ with $bd=h$ and
$a+c=t$, $ac+b+d=s$, $ad+bc=r$ for quadratic factors.
These tests are exhaustive by Gauss's lemma.

Set $d_*=A$ if $A>0$, and $d_*=B$ otherwise. For an irreducible
polynomial this is positive and nonsquare. Its pair field is
$\Q(\sqrt{d_*})$. Normality is equivalent to $J$ being a square
in this field, hence to $J$ or $Jd_*$ being an integer square. The first
case gives $V_4$, the second $C_4$. Applying these square and divisibility
tests in the displayed bounded ranges gives exactly the table.
This is an arithmetic case analysis of bounded integers, independent
of fusion-ring data.

For the scalar assertion put $A_0=1-s+h$, $B_0=t-r$,
$C_0=3t-r$, and $F_0=2s-4$. Logarithmic differentiation of the
polynomial at $\mathrm i$ gives
\[
 \operatorname{Tr}(g(\alpha))=
 4-\frac{(4C_0+3F_0)A_0+(4F_0-3C_0)B_0}{A_0^2+B_0^2}.
\]
In every row except four this is at most $583/40$, and
$g(\alpha_0)>4$. The exceptions have the following rational data:
\[
\begin{array}{c|c|c|c|c}
(t,s,r,h)&\operatorname{Tr}g&a&f(a)&f(3/10)\\\hline
(7,14,8,1)&425/29&1/10&3331/10000&-3209/10000\\
(8,14,7,1)&443/29&1/4&1/256&-479/10000\\
(11,21,11,1)&278/19&1/10&991/10000&-6989/10000\\
(12,19,8,1)&930/61&1/4&1/256&-59/10000.
\end{array}
\]
In the first and third rows $f'<0$ on $[0,1/10]$ follows from
$4x^3+2sx-r<0$. In the other two, $f''>0$ on $[0,1/4]$ and
$f'(1/4)=-23/16,-11/16$. Thus $a<\alpha_0<3/10$.
As $g$ increases up to $1/2$, subtracting $441/101$ or $81/17$
from the listed trace proves~\eqref{eq:low-trace-scalar}.
For the other rows the least root is below one: this is immediate at
norm one, while the norm-four polynomials are negative at $1/2$.

In degree two the quadratic norm is $1$ or $2$ and the trace at most
six. The large root $b$ exceeds $13/5$, and $g(a)\le5$. Hence
\[
 g(a)+2g(b)<5+2g(13/5)=\frac{1014}{97}<\frac{423}{40}.
\]
\end{proof}


\subsection{Proof of the uniform bound}

\begin{proof}[Proof of Theorem~\ref{thm:quartic-uniform-bound} outside the remaining field]
The reductions above give $u+v\le4/5$ and $h\in\{1,4\}$.
Lemma~\ref{lem:three-largest-product} gives $q_0<1/168$; hence
Lemma~\ref{lem:relative-trace-950} gives
$\operatorname{Tr}(\alpha)<14$. Its integral trace is at most
thirteen. A rational $\alpha$ is a positive integer with
$\alpha_0\le1$, hence equals one. Suppose $\alpha\ne1$.
For $i=1,2,3$, completing a square
and using $q_i=p_i/\alpha_i$ gives
\[
 p_i^2+q_i^2\ge\frac23p_i+\frac43q_i
 -\frac19\frac{(\alpha_i+2)^2}{\alpha_i^2+1}.
\]
Lemma~\ref{lem:abelian-low-trace}, including its quadratic case, yields
\[
 u+v>\frac{33}{40}-\frac23d-\frac43q_0
 >\frac45+\frac{3917}{239400}>\frac45,
\]
contradicting~\eqref{eq:quartic-two-moments}.
Consequently $E=D$; spectral saturation gives a spherical section
of the sphericalization, proving the theorem.
\end{proof}


\subsection{The remaining biquadratic field}

It remains to prove Theorem~\ref{thm:quartic-uniform-bound} for
$K=\Q(\sqrt2,\sqrt5)$, under the same assumption $D\ge950$.


We retain $p_i=1/D_i$, $q_i=1/E_i$, $d=p_0=1/D$, and
$\alpha=E/D$, choosing $E$ to be its largest signed codegree.
Thus $d\le q_0$, $q_0=\min q_i$, and $\alpha_0=d/q_0\le1$
is the least conjugate of $\alpha$. In this field the averaged
indicator inequality is
\begin{equation}\label{eq:fixed-grid}
 5u+4v\le\frac{15}{4},\qquad
 u=\sum_{i=0}^3p_i^2,\quad v=\sum_{i=0}^3q_i^2.
\end{equation}
The norm $N(\alpha)$ is a positive integer square. In particular,
\begin{equation}\label{eq:fixed-product}
 q_0^3(1-3q_0)\le\prod_iq_i\le\prod_ip_i
 \le\frac{d(1-d)^3}{27}.
\end{equation}

\begin{lemma}\label{lem:fixed-relative-pairs}
If $D\ge950$, then $q_0<1/85$ and $q_i>1/16$ for
$i=1,2,3$. Moreover
\[
 E_j+E_k\ge D_j+D_k\qquad(1\le j<k\le3).
\]
\end{lemma}
\begin{proof}
Every conjugate of $\alpha$ belongs to the signed order $T$:
\[
 \sigma(\alpha)=(E/\sigma E)^{-1}(E/\sigma D).
\]
Both factors on the right belong to $T$ by the central-column
identities, and the first denominator is a unit, since its norm is one.
For a nonidentity involution $\sigma_i$ of $K$,
$\alpha\sigma_i(\alpha)$ is therefore totally positive and belongs to
$T$. Its coefficient at the identity is a positive integer, so
\[
 I_i=dp_i(1/q_0+1/q_i)+p_jp_k(1/q_j+1/q_k)\ge1,
 \qquad\{i,j,k\}=\{1,2,3\}.
\]
Writing $U=dp_i$ and $V=p_jp_k$, one obtains
\begin{equation}\label{eq:fixed-relative-trace}
 I_i-1=U\bigl(\operatorname{Tr}E-\operatorname{Tr}D\bigr)
 +(V-U)\bigl((E_j+E_k)-(D_j+D_k)\bigr).
\end{equation}
We first prove $V>U$. Otherwise $\prod p_l=UV\le U^2<d^2$.
The first inequality of~\eqref{eq:fixed-product} then gives $q_0<1/90$,
since
\[
 \frac1{90^3}\left(1-\frac3{90}\right)>\frac1{950^2}.
\]
Consequently~\eqref{eq:fixed-grid} and
$v\ge q_0^2+(1-q_0)^2/3$ give $u<49/100$. But minimizing the square
sum at fixed $p_jp_k$ yields
\[
 p_jp_k\ge\frac12\left(\frac{(1-d)^2}{2}-u+d^2\right)
 >\frac1{200}-\frac d2>d>dp_i,
\]
a contradiction. Since $\operatorname{Tr}E\le\operatorname{Tr}D$,
equation~\eqref{eq:fixed-relative-trace} proves all three pair
inequalities.

Relabel so that $q_1=\min(q_1,q_2,q_3)$ and put $S=q_0+q_1$.
The pair inequality for $2,3$ implies
\[
 q_2q_3\le\frac{(p_2+p_3)(q_2+q_3)}4.
\]
Writing $A=1-d$, $B=1-S$, and $x=p_1$, we have
\[
 u\ge d^2+x^2+\frac{(A-x)^2}{2},\qquad
 v\ge q_0^2+q_1^2+B^2-\frac{(A-x)B}{2}.
\]
Minimization in $x$ gives
\begin{equation}\label{eq:fixed-bootstrap}
 5u+4v\ge\frac{21}{5}-\frac{32S}{5}+\frac{118S^2}{15}
 +\frac{2d}{3}(-2S+10d-3)-8q_0q_1.
\end{equation}
If $S\le3/40$, using $q_0q_1\le S^2/4$ in this inequality gives
\[
 5u+4v\ge\frac{21}{5}-\frac{32S}{5}+\frac{88S^2}{15}
 -\frac{2d}{3}(2S+3)
 \ge\frac{15}{4}+\frac3{3800},
\]
a contradiction. Thus $S>3/40$.

Equation~\eqref{eq:fixed-product} gives $q_0<1/27$.
Suppose $q_0\ge1/80$. Since $q_i>3/40-q_0$ for $i\ne0$,
\[
 \prod_iq_i>
 \psi(q_0),\qquad
 \psi(x)=x(3/40-x)^2(17/20+x).
\]
On $[1/80,1/27]$ one has
$\psi''(x)=12x^2+21x/5-39/160<0$, and
$\psi(1/27)>\psi(1/80)$. Hence
\[
 \prod_iq_i>\frac{69}{1638400}>\frac1{27\cdot950},
\]
contrary to~\eqref{eq:fixed-product}. Therefore $q_0<1/80$ and
$q_i>3/40-1/80=1/16$ for $i\ne0$. A final use of the product bound gives
\[
 q_0<\frac{256d}{27(1-1/80-1/8)}
 =\frac{20480d}{1863}\le\frac{20480}{1863\cdot950}<\frac1{85}.
\]
\end{proof}

\begin{lemma}\label{lem:fixed-small-trace}
Under the hypotheses of Lemma~\ref{lem:fixed-relative-pairs},
$\operatorname{Tr}_{K/\Q}(\alpha)<12$.
\end{lemma}
\begin{proof}
Normalize the last three coordinates by
$P_i=p_i/(1-d)$ and $Q_i=q_i/(1-q_0)$, and put
$U=\sum P_i^2$, $V=\sum Q_i^2$, and $r=\sum Q_i^{-1}$.
The preceding lemma and~\eqref{eq:fixed-grid} give
\begin{equation}\label{eq:fixed-joint-budget}
 5U+4V<\frac{15}{4}\left(\frac{85}{84}\right)^2
 <\frac{96}{25},\qquad V<\frac{163}{300},\qquad Q_i>\frac1{16}.
\end{equation}
Cauchy--Schwarz in the plane of coordinate sum zero gives
\begin{equation}\label{eq:fixed-joint-cauchy}
 \sum_i\frac{P_i}{Q_i}
 \le\frac r3+\sqrt{\left(U-\frac13\right)
              \left(\frac{2r^2}{3}-\frac{4r}{1-V}\right)}.
\end{equation}
We bound its right-hand side by $163/15$, retaining the tradeoff
between $U$ and $V$ in~\eqref{eq:fixed-joint-budget}.

For $V<1/2$, Schur's inequality gives
$Q_1Q_2Q_3\ge(1-2V)/9$, hence
\[
 r\le\frac{9(1-V)}{2(1-2V)}.
\]
Thus $r\le27/2$ if $V\le2/5$, and $r\le81/4$ if
$V\le7/16$.
For the remaining interval we use $r<108/5$. To see this, put
$t=\min Q_i>1/16$. If $t\ge1/8$, convexity gives
$r\le2/t+1/(1-2t)\le52/3$. Otherwise,
\[
 r<\frac1t+\frac{2(1-t)}{1-2t+2t^2-163/300}=:R(t).
\]
On $[1/16,1/8]$ its denominator is at least $571/2400$, so
\[
 R'(t)\le-64+
 \frac{463/150}{(571/2400)^2}<0.
\]
Therefore $r<R(1/16)=70144/3259<108/5$.

The second factor under the square root
in~\eqref{eq:fixed-joint-cauchy} increases with $r$ on the domain
$9\le r$, $1/3\le V<163/300$, and decreases with $V$.
The three intervals just considered therefore give the bounds
\[
 \sum_i\frac{P_i}{Q_i}<
 \begin{cases}
 \displaystyle\frac92+\sqrt{\frac{21}{125}\frac{81}{2}},
     &V\le2/5,\\[4pt]
 \displaystyle\frac{27}{4}+\sqrt{\frac{43}{375}\frac{1107}{8}},
     &2/5\le V\le7/16,\\[4pt]
 \displaystyle\frac{36}{5}+\sqrt{\frac{127}{1500}\frac{3936}{25}},
     &V\ge7/16.
 \end{cases}
\]
Each displayed bound is less than $163/15$, as follows by squaring.
Since $\alpha_0=d/q_0\le1$, we conclude
\[
 \operatorname{Tr}\alpha
 =\frac d{q_0}+\frac{1-d}{1-q_0}\sum_i\frac{P_i}{Q_i}
 <1+\frac{85}{84}\frac{163}{15}
 =\frac{3023}{252}<12.
\]
\end{proof}


\begin{lemma}\label{lem:fixed-arithmetic}
Let $\alpha$ be a totally positive integral $d$-number in
$K=\Q(\sqrt2,\sqrt5)$, with square norm, trace less than
twelve, and least conjugate at most one. Then $\alpha=1$, or
$\alpha$ is a conjugate of $(3+\sqrt5)/2$ or $2+\sqrt2$.
\end{lemma}
\begin{proof}
Put $\phi=(1+\sqrt5)/2$. The integer ring has basis
$1,\sqrt2,\phi,\phi\sqrt2$, since the quadratic discriminants
$8$ and $5$ are coprime. In particular its traces are even, so
$\operatorname{Tr}\alpha\le10$. The arithmetic--geometric mean
inequality gives $N(\alpha)<40$; the possible square norms are
$1,4,9,16,25,36$.

The principal ideal of a $d$-number is Galois invariant. At an
unramified prime its norm valuation is a multiple of four, since
the prime-ideal valuations above it are all equal. Only $2$ and $5$
ramify in $K$; this excludes norms $9$ and $36$. The unique prime
$\mathfrak p$ above $2$ has norm $4$ and
$\mathfrak p^2=(2)$. Thus norm $16$ would imply $\alpha=2\epsilon$
for a totally positive unit $\epsilon$. Its trace is even and at
most five, hence at most four. Equality in the arithmetic--geometric
mean inequality forces $\epsilon=1$, contrary to the least-conjugate
hypothesis.

If the norm is $25$, the $d$-number criterion forces the trace to be
divisible by five, hence to equal ten. Write
\[
 2\alpha=5+b\sqrt2+c\sqrt5+e\sqrt{10},\qquad
 c\equiv1\pmod2,\quad b\equiv e\pmod2.
\]
Positivity gives $|b|\sqrt2<5$, $|c|\sqrt5<5$ and
$|e|\sqrt{10}<5$. Thus $c=\pm1$; either $e=0$ and
$b\in\{0,\pm2\}$, or $e=\pm1$ and $b\in\{\pm1,\pm3\}$.
The norm equation is
\[
 (25+5c^2-2b^2-10e^2)^2-5(10c-4be)^2=400.
\]
For $e=0$ it forces $b=0$. For $|e|=|b|=1$ its left side is
$144$ or $-656$; for $|e|=1$, $|b|=3$ it is $-16$ or $-2416$.
The only remaining elements are $(5\pm\sqrt5)/2$, both greater
than one. Norm $25$ is therefore impossible.

The norm is now $1$ or $4$. A quartic $\alpha$ must appear in the
first seven rows of Lemma~\ref{lem:abelian-low-trace}. Their cyclic
fields are not $K$, and their two biquadratic fields are
$\Q(\sqrt2,\sqrt3)$ and $\Q(\sqrt3,\sqrt5)$.
Thus $\alpha$ has degree at most two. In degree two its quadratic
norm is one or two and its quadratic trace is at most five. At norm
one the nonrational possibilities have traces $3,4,5$ and quadratic
fields $\Q(\sqrt5)$, $\Q(\sqrt3)$,
$\Q(\sqrt{21})$, respectively. Only the first lies in $K$.
At norm two the $d$-number condition requires even trace, leaving
$2\pm\sqrt2$. A rational $\alpha$ is a positive integer with least
conjugate at most one, so equals one.
\end{proof}

\begin{proof}[Proof of Theorem~\ref{thm:quartic-uniform-bound} in the remaining field]
The preceding lemmas leave only two nonconstant dimension ratios.
Suppose first that the least conjugate is $a=(3-\sqrt5)/2$.
Since the conjugates $a,a^{-1}$ each occur twice and
$\sum p_i=\sum p_i/\alpha_i=1$, the ordinary reciprocal mass at
the two $a$ positions is $a/(1+a)$, and the signed mass is
$1/(1+a)$. One of these positions has index zero, with $p_0=d$.
Minimizing the square sums in each pair gives
\[
 5u+4v\ge\frac{7a^2+13/2}{(1+a)^2}
 -\frac{10a+8/a}{1+a}\,d
 >\frac{39}{10}-28d>\frac{15}{4}.
\]
Indeed, $a^2=3a-1$ makes the constant $21/5-1/(10a)>39/10$,
while $1/3<a<2/5$ bounds the coefficient by $28$. The last strict
inequality holds already for $D\ge256$.

If the least conjugate is $2-\sqrt2$, the ordinary mass of its two
positions is $1/2$, while their signed mass is
$(2+\sqrt2)/4>4/5$. The same pair minimization gives
\[
 u\ge\frac38-d,\qquad v>\frac{16}{25}-4d,
 \qquad 5u+4v>\frac{887}{200}-21d>\frac{15}{4}
 \quad(D\ge256).
\]
Both cases contradict~\eqref{eq:fixed-grid}. Therefore
$\alpha=E/D=1$. Spectral saturation gives a spherical section of
the sphericalization, proving the theorem.
\end{proof}

\begin{corollary}\label{cor:quartic-multiplicity-bound}
If a nonspherical self-dual rank-four complex fusion category has
quartic Frobenius--Perron global dimension, every fusion multiplicity
is at most $30$.
\end{corollary}
\begin{proof}
Theorem~\ref{thm:quartic-uniform-bound} gives $D<950<31^2$.
Lemma~\ref{core:lem:coefficient-bound} gives
$N_{ij}^{k}<\sqrt D<31$; use integrality.
\end{proof}


